\section{Separation from an ideal secret bit}
\label{sec:secret-key}

We now combine the common factors at a complete transcript with the
four-effect inequality.  A nearly correct bit has little weight in the
$01$ and $10$ branches.  A nearly uniform secret bit requires Eve's
$00$ and $11$ blocks to be close and to carry nearly half the total
weight each.  Root fidelity makes these two requirements comparable,
and the resulting estimates give a uniform numerical gap.

Throughout the paper the trace norm is unhalved:
\[
 \norm{H}_1=\tr\sqrt{H^*H}.
\]
For a classical standard Borel register with a finite-dimensional
quantum register, a state is a positive matrix-valued measure of total
trace one.  If a Hermitian difference of such measures has density
$H(t)$ with respect to a common finite dominating measure $\lambda$,
its trace norm is $\int\norm{H(t)}_1\,d\lambda(t)$.
This value is independent of the chosen dominating measure by
homogeneity.  The usual trace distance between normalized states is
one half of this norm.

We use trace-norm contraction on Hermitian differences under channels,
including classical readouts.  In finite dimensions it follows directly
from duality: the adjoint of a trace-preserving positive map is unital
and positive, so it sends each Hermitian contraction $Q$, with
$-I\preceq Q\preceq I$, to another Hermitian contraction.  Taking the
supremum of $\tr(QH)$ over such $Q$ proves contraction.  The same
argument applies to trace-class quantum registers and classical
matrix-valued measures through their bounded observables.

\begin{definition}[Distillable secret key]\label{def:distillable-key}
Let $\omega$ be a bipartite state, with Eve initially holding a
purification.  A key-distillation sequence uses an admitted protocol
of Definition~\ref{def:protocol} on each $\omega^{\otimes n}$, allowing
each party to process all its $n$ copies jointly.  The normalized
classical output at each party is a string of a fixed length
$\ell_n\in\mathbb Z_{\ge0}$ for that $n$.  Eve's output register $E'$
retains the initial purification and the complete public transcript.
For any normalized state or classical-quantum law $\sigma_{E'}$, define
the ideal length-$\ell$ key by
\[
 \gamma^{(\ell)}_\sigma
 =2^{-\ell}\sum_{z\in\{0,1\}^{\ell}}
     |z,z\rangle\langle z,z|\otimes\sigma_{E'}.
\]
The sequence is secure if there are normalized comparison laws
$\sigma_{E'}^{(n)}$ such that
\begin{equation}\label{eq:key-definition}
 \norm{\tau^{(n)}_{K_AK_BE'}-
        \gamma^{(\ell_n)}_{\sigma^{(n)}}}_1\longrightarrow0.
\end{equation}
Its asymptotic lower rate is $\liminf_{n\to\infty}\ell_n/n$.
The distillable secret key $K_D(\omega)$ is the supremum of these
rates over secure sequences, including the zero rate obtained by
producing the empty key.  A positive rate therefore requires
$\ell_n\ge1$ for all sufficiently large $n$.
\end{definition}

For positive matrices define the unnormalized root fidelity by
\[
 f(H,J)=\norm{\sqrt H\sqrt J}_1.
\]
In particular $f(cH,cJ)=c f(H,J)$ for $c\ge0$.  The next facts hold
at singular matrices as well as invertible ones.  The lower comparison
is the unnormalized form associated with the Powers--St\o rmer
inequality \cite[Lemma~4.1]{PowersStormer1970}; for normalized states
it is the lower Fuchs--van de Graaf bound
\cite[Theorem~1]{FuchsVdG1999}.  We give the short matrix proof needed
at our exact normalization.

\begin{lemma}[Root fidelity and trace norm]\label{lem:fidelity}
For arbitrary square complex matrices $A,B$ of the same size,
\begin{equation}\label{eq:polar-fidelity}
 f(A^*A,B^*B)=\norm{AB^*}_1.
\end{equation}
For positive matrices $H,J$ of the same size,
\begin{equation}\label{eq:fidelity-lower}
 f(H,J)\ge
 \frac{\tr H+\tr J-\norm{H-J}_1}{2}.
\end{equation}
Also $f(H^\T,J^\T)=f(H,J)$.
\end{lemma}

\begin{proof}
Extend the partial isometries in the polar decompositions of the square
matrices to unitaries, writing
$A=U\sqrt{A^*A}$ and $B=V\sqrt{B^*B}$.  Then
$AB^*=U\sqrt{A^*A}\sqrt{B^*B}V^*$, and unitary invariance of the trace
norm proves \eqref{eq:polar-fidelity}.  This extension is available
even when either matrix is singular.

For \eqref{eq:fidelity-lower}, put $A=\sqrt H$, $B=\sqrt J$, and
$Z=A-B$.  The matrix $Z$ is Hermitian and
$H-J=ZA+BZ$.  The spectral sign $\operatorname{sgn}(Z)$ is a Hermitian
contraction.  Trace-norm duality and cyclicity therefore give
\[
 \norm{H-J}_1
 \ge\tr[\operatorname{sgn}(Z)(H-J)]
 =\tr[|Z|(A+B)].
\]
In a unit eigenvector of $Z$ with eigenvalue $z$, positivity of $A,B$
implies
\[
 \langle A+B\rangle\ge|\langle A-B\rangle|=|z|.
\]
Summing in an eigenbasis of $Z$ shows
$\tr[|Z|(A+B)]\ge\tr Z^2$.  Finally,
\[
 \tr Z^2=\tr H+\tr J-2\tr(AB),
 \qquad 0\le\tr(AB)\le\norm{AB}_1=f(H,J).
\]
Combining these inequalities proves the lower bound.  No commutation
of $H$ and $J$ was used.  For the transpose identity, functional
calculus gives $\sqrt{H^\T}=(\sqrt H)^\T$.  The product of the two
transposed square roots is $(\sqrt J\sqrt H)^\T$; transpose and adjoint
preserve trace norm, proving $f(H^\T,J^\T)=f(H,J)$.
\end{proof}

\begin{theorem}[Uniform obstruction to a secret bit]\label{thm:zero-key}
Let $\omega\in\mathscr C$.  For every integer $n\ge1$, let an admitted
completed protocol on $\omega^{\otimes n}$ produce normalized bit
registers $K_A,K_B$, and let its output $\tau_{K_AK_BE'}$ retain in
$E'$ an initial purification held by Eve and the complete public
transcript.  The purification may be arbitrary.  For every normalized
state or classical-quantum law $\sigma_{E'}$,
\begin{equation}\label{eq:uniform-secret-bit-gap}
 \left\lVert\tau_{K_AK_BE'}-
  \frac12\sum_{i=0}^1|i,i\rangle\langle i,i|\otimes\sigma_{E'}
 \right\rVert_1\ge\frac15.
\end{equation}
Thus the conventional trace distance is at least $1/10$, and
$K_D(\omega)=0$ in Definition~\ref{def:distillable-key}.
\end{theorem}

\begin{proof}
Group all of Alice's factors and all of Bob's factors, and set
$R=\omega^{\otimes n}$.  Lemma~\ref{lem:tensor-closure} gives
$R\in\mathscr C$.  First apply the classical readout of Eve's retained
complete public transcript, leaving her initial purifying register
unchanged, and discard any extra retained registers.  On the actual output
this produces the full classical transcript together with the purification.
On any ideal bit comparison state it produces an ideal bit with a
classical-quantum comparison law.  Next,
Lemma~\ref{lem:purification-recovery} supplies a channel from the retained
purifying space to the canonical one, leaving the transcript unchanged.
It recovers the actual canonical output and again sends every ideal law
to an ideal law.  Trace-norm contraction under these Eve-side channels
therefore reduces the claim to the canonical purification
\eqref{eq:canonical-purification} and the full classical transcript $T$.

Apply Lemma~\ref{lem:transcript} to obtain a finite transcript measure
$\mu$, the positive factors $X_i(t),Y_j(t)$, and the Eve densities
$\tau_{ij}(t)$ of \eqref{eq:eve-blocks}.  Fix any normalized positive
matrix-valued comparison law $\sigma_{ET}$ on the same transcript space.

Let $\nu$ be the transcript law of $\sigma_{ET}$, namely
$\nu(S)=\tr\sigma_{ET}(S)$, and use the common finite measure
$\lambda=\mu+\nu$.  If $h=d\mu/d\lambda$, replace each actual density
by $h(t)\tau_{ij}(t)$ and both $X_i(t)$ by $h(t)X_i(t)$, leaving the
$Y_j(t)$ unchanged.  Equation~\eqref{eq:eve-blocks} then still holds,
including where $h=0$.  The ideal law has a positive density
$\sigma(t)$ with respect to $\lambda$ and
$\int\tr\sigma(t)\,d\lambda(t)=1$.  Homogeneity preserves all actual
probabilities and root-fidelity integrals in this change of measure.
We now denote $\lambda$ again by $\mu$.  This common measure keeps
singular parts of either original transcript law; there is no
assumption that the ideal and actual laws have the same support.

For these common densities set
\[
 D_{ij}(t)=X_i(t)\otimes Y_j(t),\qquad
 p_{ij}(t)=\tr\tau_{ij}(t)=\tr[D_{ij}(t)R],\qquad
 P_{ij}=\int p_{ij}(t)\,d\mu(t).
\]
All $P_{ij}$ are nonnegative and $\sum_{i,j}P_{ij}=1$.  With
$A=\sqrt{D_{00}(t)}\sqrt R$ and
$B=\sqrt{D_{11}(t)}\sqrt R$, the transpose invariance and
\eqref{eq:polar-fidelity} give
\begin{equation}\label{eq:transcript-fidelity}
 \begin{aligned}
 f(\tau_{00}(t),\tau_{11}(t))
 &=f\bigl(\sqrt R D_{00}(t)\sqrt R,
           \sqrt R D_{11}(t)\sqrt R\bigr)\\
 &=\norm{\sqrt{D_{00}(t)}R\sqrt{D_{11}(t)}}_1.
 \end{aligned}
\end{equation}
The privacy interpretation of an off-diagonal coherence norm as the
fidelity of conditional Eve states is also used in the private-state
approach \cite[Eq.~(7)]{keyBoundEntanglement}.  Here the identity is
applied to unnormalized branches at each complete transcript.  Its
proof includes singular $R$ and singular effects.

Apply Lemma~\ref{lem:coherence} at almost every $t$.  Scalar
Cauchy--Schwarz after integration gives
\begin{equation}\label{eq:integrated-coherence}
 \begin{split}
 \mathfrak f
 &:=\int f(\tau_{00}(t),\tau_{11}(t))\,d\mu(t)\\
 &\le\sqrt{P_{00}P_{01}}+\sqrt{P_{10}P_{11}}
       +2\sqrt{P_{01}P_{10}}.
 \end{split}
\end{equation}
For example, $\int\sqrt{p_{00}(t)p_{01}(t)}\,d\mu(t)$ is at most
$\sqrt{P_{00}P_{01}}$.  The factors need only be positive; no
pointwise POVM normalization was used.

We next express what small distance from the fixed ideal bit would
require of this same $\mathfrak f$.  Let $\eta$ be the unhalved
trace-norm distance after the common transcript readout, and define
\[
 \begin{split}
 e&=P_{01}+P_{10},\\
 q&=\int\norm{\tau_{00}(t)-\tau_{11}(t)}_1\,d\mu(t),\\
 a_i&=\int\norm{\tau_{ii}(t)-\tfrac12\sigma(t)}_1\,d\mu(t)
       \quad(i=0,1).
 \end{split}
\]
The unequal-bit blocks of the ideal state are zero.  Additivity of the
trace norm across the classical bit blocks, and then the triangle
inequality within the two equal-bit blocks, give
\begin{equation}\label{eq:linked-errors}
 \eta=e+a_0+a_1,\qquad q\le a_0+a_1=\eta-e.
\end{equation}
Trace conservation gives one further inequality that is needed for the
uniform constant.  The actual equal-bit blocks have total trace $1-e$,
whereas the corresponding ideal blocks have total trace one.  Thus
\[
 a_0+a_1
 \ge\left|\sum_{i=0}^1\int
       \tr\bigl(\tau_{ii}(t)-\tfrac12\sigma(t)\bigr)\,d\mu(t)\right|
 =e,
\]
and \eqref{eq:linked-errors} implies
\begin{equation}\label{eq:unequal-error-bound}
 e\le\frac\eta2.
\end{equation}

Integrating \eqref{eq:fidelity-lower} gives
\[
 \mathfrak f\ge\frac{P_{00}+P_{11}-q}{2}
 =\frac{1-e-q}{2}\ge\frac{1-\eta}{2}.
\]
On the other hand, Cauchy--Schwarz on the first two terms of
\eqref{eq:integrated-coherence} and the arithmetic--geometric mean
inequality on the last give
\[
 \begin{aligned}
 \sqrt{P_{00}P_{01}}+\sqrt{P_{10}P_{11}}
 &\le\sqrt{(P_{00}+P_{11})(P_{01}+P_{10})}
   =\sqrt{e(1-e)},\\
 2\sqrt{P_{01}P_{10}}&\le e.
 \end{aligned}
\]
The two estimates have the exact combined form
\begin{equation}\label{eq:fidelity-sandwich}
 \frac{1-\eta}{2}\le\mathfrak f\le\sqrt{e(1-e)}+e,
 \qquad e\le\frac\eta2.
\end{equation}
If $\eta<1/5$, then $e<1/10$.  On $0\le e\le1/10$ the function
$\sqrt{e(1-e)}+e$ is increasing and has value $2/5$ at $e=1/10$.
The right side of \eqref{eq:fidelity-sandwich} is therefore strictly
less than $2/5$, while its left side is strictly greater than $2/5$.
This contradiction proves \eqref{eq:uniform-secret-bit-gap},
uniformly in $n$, the represented state, the protocol, and the ideal
comparison law.

Finally suppose a secure key-distillation sequence had $\ell_n\ge1$
for infinitely many $n$.  Along those indices, the local classical channels
that retain the first bit of each string and discard the other bits,
while retaining $E'$, map $\gamma^{(\ell_n)}_{\sigma^{(n)}}$ to an ideal
bit: exactly half the uniform strings have either first bit.
Trace-norm contraction would make the resulting bit-output error tend to
zero, contradicting \eqref{eq:uniform-secret-bit-gap}.  Thus every secure
sequence has $\ell_n=0$ for all sufficiently large $n$.  In particular its
asymptotic lower rate is zero.  The empty key achieves rate zero, so
$K_D(\omega)=0$.
\end{proof}

The proof also yields the retained conditioning and limiting forms of
the obstruction.  These statements specify how normalization and
Eve's complete information are preserved.

\begin{proposition}[Conditioning, aborts, and completed-output limits]
\label{prop:operational-extensions}
Consider inputs $\omega^{\otimes n}$ with $\omega\in\mathscr C$ and
$n\ge1$, and admitted protocols whose normalized completed outputs retain
in Eve's register an arbitrary initial purification and the complete
public transcript.  The following statements extend the gap in
Theorem~\ref{thm:zero-key}; in item 2, the completed local output alphabet
may include an abort symbol.
\begin{enumerate}
\item If $A$ is a measurable event of the complete public transcript
with probability $r>0$, the normalized bit output conditioned on $A$
has distance at least $1/5$ from every ideal bit law for that accepted
output.  This holds for every $r>0$.
\item Suppose a sequence of completed protocols has a locally
available abort symbol at either party.  If the probability that at
least one party aborts tends to zero, and the normalized non-aborting
outputs approach ideal non-aborting bit laws, then replacing each
local abort by the fixed bit zero would give completed bit outputs
approaching ideal bits.  Such a sequence is impossible for inputs in
$\mathscr C$ (including any numbers of copies).
\item Let $\tau_k$ be completed normalized bit outputs to which the
gap applies, possibly on different output spaces.  Suppose their bit
registers are kept fixed and their Eve registers are embedded in a
common output space $F$, giving states $\widehat\tau_k$.  The embeddings
must preserve and permit recovery of the full initial purification and
complete public records, without coarse-graining those records.  In
particular, an Eve-side recovery channel recovers $\tau_k$ from
$\widehat\tau_k$ while leaving the bits fixed.  If
$\widehat\tau_k\to\tau$ in trace norm, then for every normalized law
$\sigma_F$,
\[
 \norm{\tau-\tfrac12\sum_{i=0}^1
       |i,i\rangle\langle i,i|\otimes\sigma_F}_1\ge\frac15.
\]
\end{enumerate}
\end{proposition}

\begin{proof}
For the first assertion, use the canonical factors from
Lemma~\ref{lem:transcript}.  Restrict their measure $\mu$ to $A$ and
replace both $X_0(t),X_1(t)$ by $X_0(t)/r,X_1(t)/r$, leaving the
$Y_j(t)$ unchanged.  Each of the four blocks in
\eqref{eq:eve-blocks} is divided by $r$, and its integral over $A$ is
the corresponding normalized accepted block.  All probabilities and
the root fidelity scale by the same factor $1/r$.  The accepted total
trace is one, so the proof of \eqref{eq:fidelity-sandwich}, including
$e\le\eta/2$, applies without change.  The recovery and common-measure
arguments cover arbitrary purifications and arbitrary ideal laws for
the accepted output.  No positive lower bound on $r$ was used.

For the second assertion, let $\delta_k$ be the probability of at
least one abort, let $\tau_k^{\mathrm{ok}}$ be the normalized
non-aborting state when $\delta_k<1$, and let $\xi_k$ be the normalized
state on the aborting part when $\delta_k>0$.  The completed state is
the sum of these two positive classical parts with weights
$1-\delta_k$ and $\delta_k$.  Let $\mathcal Q$ be the product of the
two local classical channels sending each party's own abort symbol to
zero and fixing its bit values.  It is defined on every completed
output and preserves each non-aborting ideal bit $\gamma_k$.  Since
the trace-norm distance between normalized states is at most two,
\[
 \norm{\mathcal Q(\tau_k)-\gamma_k}_1
 \le (1-\delta_k)\norm{\tau_k^{\mathrm{ok}}-\gamma_k}_1
       +2\delta_k.
\]
The right side tends to zero under the stated hypotheses.  This
contradicts Theorem~\ref{thm:zero-key} for the completed bit-output
protocol with the local replacement appended.  The replacement uses
only abort symbols actually available to the respective laboratories.

For the third assertion, denote the recovery channel for the $k$th
embedding by $\mathcal R_k$.  For any ideal bit $\gamma_\sigma$ in the
common space, applying this channel to Eve maps it to the ideal bit
$\gamma_{\mathcal R_k(\sigma)}$ on the $k$th output space and recovers
$\tau_k$ from $\widehat\tau_k$.  Contraction and the per-protocol gap
therefore give
\[
 \norm{\widehat\tau_k-\gamma_\sigma}_1
 \ge\norm{\tau_k-\gamma_{\mathcal R_k(\sigma)}}_1
 \ge\frac15.
\]
The triangle inequality now yields
\[
 \norm{\tau-\gamma_\sigma}_1
 \ge\frac15-\norm{\widehat\tau_k-\tau}_1
 \longrightarrow\frac15.
\]
Only the completed actual outputs need converge.  The embeddings must
retain and permit recovery of the full purification and public records,
so they do not coarse-grain that information.  The argument requires
no convergence on the auxiliary probe inputs and no factorization
statement for the limit itself.  In particular, it does not assume
that finite padded transcript laws converge in trace norm to a law of
an infinite transcript.
\end{proof}
